%! TEX root = ../main.tex
This report has investigated three interrelated dimensions of economic inequality and intergenerational transmission in Chile, combining microeconometric estimation, survey microdata, administrative tax records, and external wealth registries. Together, the findings provide a comprehensive empirical assessment of persistence and concentration while highlighting the critical role of modern econometric corrections.

\subsection*{Summary of Empirical Findings}

First, our replication and extension of \citet{nunez2010mobility} in Section~\ref{sec:intergenerational-mobility} demonstrates that Chile exhibits substantial intergenerational income persistence. The estimated TSTSLS income elasticities range between 0.61 (conditioning on paternal schooling and occupation) and 0.83 (schooling only), confirming that parental background exerts a first-order influence on offspring earnings. Beyond the baseline replication, our extensions reveal that labor informality introduces significant heterogeneity into the transmission mechanism: sons contributing to the pension system display substantially higher rank--rank persistence (slope of 0.320) than informal workers (0.295; Spearman rank correlation of 0.250). Furthermore, our multiple imputation sensitivity analysis underscores that the assumed correlation between paternal and filial unobservables critically shapes estimated persistence, demonstrating that social security affiliation is an essential institutional mediator of mobility in developing economies.

Second, our analysis of household wealth inequality in Section~\ref{sec:wealth-inequality} highlights the pervasive limitations of self-reported survey data. Standard surveys like the EFH suffer from unit non-response at the extreme right tail, and the presence of negative net worth ($L_i > A_i$) distorts synthetic inequality metrics such as the Gini coefficient, mathematically breaking its standard axiomatic bounds (Appendix~\ref{app:gini-negative}). By implementing the semi-parametric Pareto tail correction of \citet{vermeulen2018} augmented with the \textit{Forbes World's Billionaires} census and benchmarking aggregates to National Accounts, we demonstrate that top wealth shares in Chile are substantially higher than raw survey tabulations indicate. Microfoundations of multiplicative accumulation and scale-dependent returns (Appendix~\ref{app:pareto-foundations}) justify the power-law tail, showing that capital concentration cannot be captured without accounting for ultra-high-net-worth individuals.

Third, in Section~\ref{sec:bfm-methodology}, we evaluated the frontier non-parametric calibration framework of \citet{blanchet2022weight} (BFM) using CASEN 2017 and 2022 microdata combined with tabulated administrative tax records from the SII. In contrast to parametric Pareto adjustments that rely on an exogenously chosen cutoff $w_{\min}$, the BFM approach endogenously locates the survey-tax merging point $y^*$ where marginal representativeness matches cumulative coverage. We provided formal analytical proofs establishing the sufficient conditions for the existence and uniqueness of this merging point in both continuous and discrete settings (Appendix~\ref{app:merging-point-proof}). Empirically, BFM calibration shifts substantial mass to the upper decile, increasing the top-1\% income share by over 4.4~pp in 2017 (from 9.17\% to 13.56\%, and 14.42\% with support expansion). For 2022, post-pandemic nominal income growth and tax bracket restructuring pushed the merging point to an extreme threshold of $y^* \approx 190.0$~M CLP/year (leaving only 12 survey observations above $y^*$), illustrating how calibration constraints bind when survey coverage is thin and proving the necessity of rank-preserving support expansion.

\subsection*{Methodological and Policy Implications}

The unified takeaway from this investigation is that no single empirical source is sufficient for inequality and mobility measurement. Household surveys remain indispensable for observing demographic attributes, informal labor, and the bottom half of the distribution, but fail systematically in the tail due to non-response and under-reporting. Conversely, administrative tax records provide comprehensive tracking of formal capital and labor at the top, but exclude informal workers and low-income exempt individuals. 

Rigorous distributional analysis and effective public finance design---such as evaluating wealth taxes, progressive income reforms, and social safety nets---require hybrid methodologies that combine administrative registries with representative surveys under explicit calibration principles. The open-source empirical toolkit developed in this project provides a modular, reproducible foundation for implementing these state-of-the-art corrections in Chile and across developing contexts.

